\documentclass[11pt]{article}
\usepackage{amsmath,amssymb,amsthm}
\usepackage[margin=1in]{geometry}

\newtheorem{definition}{Definition}
\newtheorem{lemma}{Lemma}
\newtheorem{proposition}{Proposition}

\newcommand{\Sv}{S_{\mathrm{v2}}}
\newcommand{\Av}{A_{\mathrm{v2}}}
\newcommand{\Incv}{\mathrm{Inc}_{\mathrm{v2}}}
\newcommand{\Aut}{\mathrm{Aut}}
\newcommand{\Rv}{R_{\mathrm{v2}}}
\newcommand{\Dv}{D_{\mathrm{v2}}}
\newcommand{\ReadoutSrc}{\mathrm{Readout}_{\mathrm{src}}}
\newcommand{\DetSrc}{\mathrm{Det}_{\mathrm{src}}}
\newcommand{\GEff}{g_{\mathrm{eff}}}

\title{Finite Toy Response v2 Source Specification}
\author{The AEther-Flow Research Project}
\date{2026-07-08}

\begin{document}
\maketitle

\section*{Control Status}

This artifact implements v18 P6-T01. It is a draft/control finite toy
response v2 source specification for \texttt{RT-20260708-014}. It is not
canonical ontology, not \(\DetSrc\) adoption, not \(\ReadoutSrc\) adoption,
not detector semantics, not source detector/readout semantics adoption, not
source-law adoption, not coupling-law adoption, not \(\GEff\), not matter
coupling, not stress-energy semantics, not a matter action, not Einstein
equations, not benchmark evidence, not a Gate Chair verdict, and not
completed-derivation evidence.

The purpose of the packet is to specify a finite source-to-response toy target
that does not depend on explicit tags whose erasure destroys the response
relation.

\section{Finite Toy Response v2 Fields}

\begin{verbatim}
finite_toy_response_v2:
  finite_source_set: "S_v2 = {a,b,c}"
  source_relation_family: "A_v2 path adjacency and Inc_v2 incidence records"
  invariant_orbit_structure: "Aut(P3) vertex and unordered-pair orbits"
  induced_response_relation: "R_v2({x,y}) = d_A(x,y)"
  metric_response_analogue: "finite graph-distance form D_v2, not g_eff"
  explicit_target_tags_forbidden: true
  target_metric_import_forbidden: true
  empirical_detector_import_forbidden: true
  relabeling_invariance_test: "distance is invariant under Aut(P3)"
  tag_removal_stress: "erase presentation labels but retain source relation"
  detector_readout_status: "SourceReadoutCandidate_EStar_v1
    candidate_placeholder_nonadopted"
\end{verbatim}

\section{Source Set and Relation Family}

\begin{definition}[Finite source set]
Let
\[
  \Sv=\{a,b,c\}.
\]
The letters \(a,b,c\) are presentation names only. They are not source tags,
target labels, detector labels, proper-time labels, or metric coordinates.
\end{definition}

\begin{definition}[Source relation family]
The source relation family consists of the path adjacency relation
\[
  \Av=\{\{a,b\},\{b,c\}\}
\]
and the incidence record
\[
  \Incv=\{(a,\{a,b\}),(b,\{a,b\}),(b,\{b,c\}),(c,\{b,c\})\}.
\]
No element of \(\Av\) or \(\Incv\) carries an explicit orientation tag,
normalization tag, target metric component, empirical detector protocol, or
benchmark-success flag.
\end{definition}

\section{Invariant Orbit Structure}

\begin{definition}[Relabeling group]
Let \(\Aut(\Sv,\Av)\) be the automorphism group of the three-vertex path.
It has two elements: the identity and the endpoint swap
\[
  \sigma(a)=c,\qquad \sigma(b)=b,\qquad \sigma(c)=a.
\]
\end{definition}

The vertex orbits are \(\{b\}\) and \(\{a,c\}\). The unordered-pair orbits
are:
\[
  \mathcal{O}_0=\{\{a,a\},\{b,b\},\{c,c\}\},
\]
\[
  \mathcal{O}_1=\{\{a,b\},\{b,c\}\},
\]
\[
  \mathcal{O}_2=\{\{a,c\}\}.
\]
These orbit classes are determined by the source relation. They are not
explicit response tags.

\section{Induced Response Relation}

\begin{definition}[Finite source distance]
Let \(d_{\Av}(x,y)\) be the shortest-path distance in the graph
\((\Sv,\Av)\). Thus
\[
  d_{\Av}(x,x)=0,\qquad
  d_{\Av}(a,b)=d_{\Av}(b,c)=1,\qquad
  d_{\Av}(a,c)=2.
\]
\end{definition}

\begin{definition}[Induced response relation]
The finite toy response relation is the source-induced partial map
\[
  \Rv:\binom{\Sv}{\leq 2}\to\{0,1,2\},\qquad
  \Rv(\{x,y\})=d_{\Av}(x,y).
\]
Equivalently, \(\Rv(\mathcal{O}_k)=k\) for \(k\in\{0,1,2\}\).
\end{definition}

\begin{definition}[Finite distance-form analogue]
The metric-response analogue is the finite graph-distance form
\[
  \Dv(x,y)=d_{\Av}(x,y).
\]
\(\Dv\) is only a finite toy distance-form analogue. It is not
\(\GEff\), not a target metric, not proper-time normalization, and not a
benchmark metric.
\end{definition}

\section{Relabeling Invariance Test}

\begin{lemma}[Relabeling invariance]
For every \(\sigma\in\Aut(\Sv,\Av)\) and every \(x,y\in\Sv\),
\[
  \Rv(\{\sigma(x),\sigma(y)\})=\Rv(\{x,y\}).
\]
\end{lemma}

\begin{proof}
Every automorphism of \((\Sv,\Av)\) preserves adjacency and therefore
shortest-path distance. Since \(\Rv\) is defined as \(d_{\Av}\), it is
constant on automorphism orbits of unordered pairs.
\end{proof}

\section{Tag-Removal Stress}

\begin{definition}[Presentation-label erasure]
Let \(E_{\mathrm{lab}}\) erase the presentation names \(a,b,c\) while
retaining the isomorphism class of the finite source relation
\((\Sv,\Av,\Incv)\). The result is the unlabeled three-vertex path.
\end{definition}

\begin{proposition}[Tag-removal survival target]
Under \(E_{\mathrm{lab}}\), the induced response relation remains defined on
the unlabeled path orbit classes. The response values are still
diagonal \(0\), adjacent \(1\), and endpoint-pair \(2\).
\end{proposition}

\begin{proof}
The construction of \(\Rv\) uses only source adjacency and shortest-path
distance. Presentation names are not inputs to \(d_{\Av}\). Erasing names
therefore removes no datum required to compute orbit distance. The unlabeled
path still has the three orbit classes diagonal, adjacent, and endpoint-pair.
\end{proof}

\section{Difference From the Frozen Explicit-Tag Route}

The older finite toy route used explicit response tags
\((\varepsilon,\lambda,\tau)\). The Refuter stress for that route showed that
erasing those tags made the response relation undefined and froze that local
route. This P6-T01 target differs in the following precise way:

\[
  \text{old route: response read from explicit tags},
\]
\[
  \text{P6-T01 route: response induced from source relation orbits}.
\]

The P6-T01 target does not claim that the finite model has been constructed
or that all later stress tests will pass. It only supplies a non-tag-fragile
source specification for P6-T02.

\section{Detector/Readout Status}

\begin{verbatim}
detector_readout_status:
  status: "candidate_placeholder_nonadopted"
  candidate: "SourceReadoutCandidate_EStar_v1"
  candidate_authority: "draft/control only"
  used_to_define_R_v2: false
  Det_src_adopted: false
  Readout_src_adopted: false
  detector_semantics_adopted: false
  matter_coupling_derived: false
\end{verbatim}

The source detector/readout candidate may be carried forward as a
nonadopted placeholder for later compatibility checks. It is not used as a
premise of \(\Rv\), and it does not provide empirical detector semantics.

\section{P6-T02 Construction Obligations}

The next packet should either construct the finite toy response v2 model from
this specification or record a precise obstruction. The construction attempt
must check:

\begin{enumerate}
  \item whether \((\Sv,\Av,\Incv,\Rv,\Dv)\) can be represented without hidden
        target metric import;
  \item whether finite relation perturbations preserve a nontrivial response
        class or produce a scoped obstruction;
  \item whether `detector_readout_status' remains placeholder-only; and
  \item whether no old explicit tag \((\varepsilon,\lambda,\tau)\) or
        equivalent tag surrogate is reintroduced.
\end{enumerate}

\section{Distance-to-GR Status}

\begin{center}
\begin{tabular}{p{0.33\linewidth}p{0.57\linewidth}}
\hline
Burden & Status \\
\hline
Source ontology primitives & Unchanged; no canonical ontology edit is made. \\
Source equivalence EqSrc & Unchanged; no general EqSrc discharge. \\
RetainH & Unchanged; not adopted. \\
GenH & Unchanged; not adopted. \\
ObsLoc\_lc & Unchanged. \\
Resp\_lc & Unchanged; this is a finite toy response target only. \\
M\_src & Unchanged; no source manifold construction is supplied. \\
\(\GEff\) & Not constructed; \(\Dv\) is a finite graph-distance analogue only. \\
matter coupling & Not derived or adopted. \\
Einstein equations & Not started. \\
finite-variation robustness & Source specification defines relabeling and tag-removal tests for later construction and stress. \\
benchmark promotion & Blocked; no benchmark evidence or promotion is supplied. \\
Gate Chair status & Unchanged; no verdict issued. \\
current route freeze or hard-fail status & Not frozen for P6-T01; construction or obstruction is routed to P6-T02. \\
\hline
\end{tabular}
\end{center}

\section{Next Route}

The next route is P6-T02:
\[
  \texttt{finite\_toy\_response\_v2\_model\_or\_obstruction}.
\]

\begin{thebibliography}{9}

\bibitem{v18-plan}
The AEther-Flow Research Project. (2026, July 7). \emph{Recommendations
implementation plan continue task v18} [Internal implementation plan].

\bibitem{handoff706}
The AEther-Flow Research Project. (2026, July 8). \emph{Handoff 0706}
[Internal research-control handoff].

\bibitem{burden-map}
The AEther-Flow Research Project. (2026, June 17). \emph{GR derivation
burden map} [Internal control note].

\bibitem{frozen-toy}
The AEther-Flow Research Project. (2026, June 18). \emph{Resp\_lc finite toy
metric-response model Refuter stress test} [Internal research-control TeX
artifact].

\end{thebibliography}

\end{document}
